\labsection{8}{Evaluate probabilistic classification}{sec:lab08}

\begin{labproject}{Lab 8 classification}
\textbf{Datasets.} \texttt{heart\_failure\_dataset.csv}, with Breast Cancer and Optical Digits extensions.\\
\textbf{Goal.} Fit a logistic model, preserve probabilistic output, and evaluate different kinds of classification error.

\textbf{Procedure.}
\begin{enumerate}
  \item Identify the binary target and convert categorical predictors to factors.
  \item Create a stratified train/test split so outcome prevalence remains comparable.
  \item Fit logistic regression using demographic, symptom, exercise, and clinical measurements.
  \item Generate holdout probabilities, then apply a 0.50 threshold to obtain class decisions.
  \item Build the confusion matrix and calculate accuracy, sensitivity, specificity, and precision.
  \item Inspect the probability distributions for the two observed classes and explain why overlap is unavoidable.
\end{enumerate}

\textbf{Interpretation checkpoint.} This exercise is about statistical classification, not clinical diagnosis. Dataset labels, measurement processes, and error costs must be understood before any real-world medical use.

\textbf{Extension comparison.} First compare \texttt{radius\_mean} and \texttt{texture\_mean} by diagnosis in the Breast Cancer dataset. Then open the Optical Digits training/test files from the supplied ZIP and reconstruct at least one $8\times8$ digit from its 64 pixel-intensity features. Identify how evaluation must change when the target has ten classes rather than two.

\lstinputlisting[style=dsr]{all_experiments/lab08_classification.R}
\end{labproject}
